% =====================================================================
\chapter{Soluciones}
\label{ap:soluciones}
% =====================================================================

Los resultados finales se dan con tres cifras significativas y los intermedios con cuatro, para que el redondeo no se acumule. Las remisiones «M24, §2.3» indican el módulo y la sección de sus apuntes, comprobadas en sus índices. Al final de cada solución figuran, entre corchetes, los criterios de corrección.

% ---------------------------------------------------------------------
\section{Parte A: clave de respuestas}

\begin{table}[h]
\centering\footnotesize
\caption{Clave de la Parte~A.}\label{tab:clave}
\begin{tabularx}{\textwidth}{@{}c c >{\raggedright\arraybackslash}X l@{}}
\toprule
N.º & Resp. & Justificación & Repasar \\
\midrule
1 & (c) & La menor cota válida es el máximo factible de $x_t$: $\min(120,\,90)=90$~t (proposición~2.11 de M24). & M24, §2.3 \\
2 & (a) & Holgura complementaria: $y_i\,(b_i-(Ax)_i)=0$ y la holgura es 12~h, luego $y_i^{\ast}=0$. & M24, §2.2 \\
3 & (d) & Fractil $30/40=0{,}75$; $q^{\ast}=200+40\cdot 0{,}6745=227{,}0$~hL. 234 usa el fractil $0{,}8$. & M24, §6.2 \\
4 & (b) & Pesos positivos dan puntos Pareto-óptimos; los no apoyados no se alcanzan (proposición~5.5). & M24, §5.4 \\
5 & (c) & $z=0$: $\EI=0\cdot\Phi(0)+2\,\varphi(0)=2\cdot 0{,}3989=0{,}798$; la EI es positiva aunque $\mu=\fmejor$. & M25, §4.3 \\
6 & (a) & $\sigma^{2}=1-k_\ast^{2}=1-(e^{-1/2})^{2}=1-e^{-1}=0{,}632$. & M25, §3.3 \\
7 & (d) & $n\ge\ln 0{,}05/\ln 0{,}95=58{,}4$, luego $n=59$. & M25, §6.2 \\
8 & (b) & Con la varianza predictiva, $\sigma$ no baja de $\sigma_\varepsilon$ en los puntos ensayados (observación~4.8). & M25, §4.5 \\
9 & (a) & Autovalor de Euler $1-5\Delta t$; $\abs{1-5\Delta t}<1\iff\Delta t<2/5=0{,}4$~h. & M26, §1.3 \\
10 & (c) & $\min_u\{u^{2}+(2+u)^{2}\}$: $2u+2(2+u)=0$, $u=-1$; $K_0=P_1/(R+P_1)=1/2$. & M26, §3.1 \\
11 & (d) & Penalización exacta: basta $\rho\ge\max_i\lambda_i^{\ast}=3$ (proposición~3.9). & M26, §3.3 \\
12 & (b) & Truncar no es terminar: $d=\texttt{terminated}=0$ y $y=r+\gamma Q_{\phi'}(s',\tilde a')$. & M26, §4.4 \\
13 & (c) & $E_e/E_0=x_e^{p}=0{,}6^{3}=0{,}216$. & M27, §2.2 \\
14 & (a) & La flexibilidad es autoadjunta: $\lambda=U$ y $\partial c/\partial x_e=-p\,x_e^{p-1}E_0\,u_e\T k_0u_e$. & M27, §2.3 \\
15 & (d) & $d^{\ast}=\varepsilon H/(2\tau\tan 30^{\circ})=0{,}02\cdot 1080/1{,}1547=18{,}7$~m. & M27, §5.3 \\
16 & (a) & Corolario~1.4: $\abs{F}=2\abs{V}-4=1996$ (y $\abs{E}=2994$). & M27, §1.2 \\
17 & (b) & $\sigma\sqrt{2\ln K}=0{,}04\sqrt{2\ln 50}=0{,}04\cdot 2{,}797=0{,}112$. & M28, §2.2 \\
18 & (d) & $0{,}1\ln 1{,}2+(-0{,}1)\ln 0{,}8=0{,}01823+0{,}02231=0{,}0405$. & M28, §4.3 \\
19 & (c) & $(K+1)/\eta^{K}=4/27=0{,}148$. & M28, §3.3 \\
20 & (b) & $a_{4;\,0{,}1}=3{,}1699$; $\VAN=20\,000\cdot 3{,}1699-50\,000=13\,397$~€. & M28, §5.1 \\
\bottomrule
\end{tabularx}
\end{table}

% ---------------------------------------------------------------------
\section{Parte B: soluciones de los problemas}

% ------------------------------------------------------------- B.1
\begin{solucion}{eje:milp}
\datos $D=100$~t; $f=(400,150,250)$~€; $c=(12,16,14)$~€/t; $L=(30,20,25)$~t; $U=(80,50,60)$~t.
\planteamiento PLEM con costes fijos y lotes mínimos (M24, §2.3); la menor $M$ válida es el máximo factible de $x_i$ (proposición~2.11 de M24); la relajación lineal acota el óptimo (proposición~2.10) y la disyunción se escribe con una binaria (M24, §2.3, apartado~(d) de la lista de modelizaciones lógicas).
\pasos

\noindent\textbf{(a)}
\begin{enumerate}[label=\arabic*.]
  \item Modelo:
  \[
  \begin{aligned}
  \min\ &\sum_{i=1}^{3}\bigl(f_iy_i+c_ix_i\bigr)\\
  \text{s.a.}\ &x_1+x_2+x_3=100,\qquad L_iy_i\le x_i\le M_iy_i,\qquad x_i\ge 0,\quad y_i\in\{0,1\}\quad(i=1,2,3).
  \end{aligned}
  \]
  \item La igualdad cubre exactamente la carga del turno. $x_i\le M_iy_i$ es el enlace: si el horno está apagado ($y_i=0$) no trata nada y, si trata algo, se paga $f_i$. $L_iy_i\le x_i$ es el lote mínimo: o $x_i=0$ o $x_i\ge L_i$.
\end{enumerate}

\noindent\textbf{(b)}
\begin{enumerate}[label=\arabic*.]
  \item En el conjunto factible sin el enlace, $x_i\le U_i$ (capacidad) y $x_i\le D=100$ (la igualdad con $x\ge 0$). Luego $\sup x_i=\min(U_i,D)$: $M=(80,\,50,\,60)$~t.
  \item Con $y_i=1$ la restricción es $x_i\le U_i$, que no corta nada factible; con $y_i=0$ fuerza $x_i=0$, como se desea. Una $M_i$ menor cortaría planes factibles; una mayor es válida pero debilita la relajación (proposición~2.11).
\end{enumerate}

\noindent\textbf{(c)}
\begin{enumerate}[label=\arabic*.]
  \item Un horno solo no basta: $U_i<100$ para todo $i$. Los patrones $(0,0,0)$, $(1,0,0)$, $(0,1,0)$ y $(0,0,1)$ son infactibles.
  \item $\{1,2\}$: $x_1+x_2=100$, $x_1\le 80$, $x_2\ge 20$. El horno~1 es más barato por tonelada, así que $x_1=80$, $x_2=20$: $550+960+320=1830$~€.
  \item $\{1,3\}$: $x_3\ge 25$ obliga a $x_1\le 75$: $x_1=75$, $x_3=25$: $650+900+350=1900$~€.
  \item $\{2,3\}$: el horno~3 es más barato: $x_3=60$, $x_2=40$: $400+840+640=1880$~€.
  \item $\{1,2,3\}$: mínimos $x_2=20$, $x_3=25$ y $x_1=55\le 80$: $800+660+320+350=2130$~€.
\end{enumerate}

\noindent\textbf{(d)}
\begin{enumerate}[label=\arabic*.]
  \item En la relajación conviene $y_i=x_i/M_i$ (el menor valor permitido; cumple $L_iy_i\le x_i$ porque $L_i\le M_i$). El problema se reduce a $\min\sum_i(c_i+f_i/M_i)x_i$ con $\sum_ix_i=100$, $0\le x_i\le M_i$.
  \item Costes efectivos: $12+400/80=17$, $16+150/50=19$ y $14+250/60=18{,}17$~€/t. Se llena el horno~1 (80~t) y después el~3 (20~t): $z_{\mathrm{LP}}=80\cdot 17+20\cdot 18{,}17=1723{,}3$~€, con $y=(1;\ 0;\ 1/3)$.
  \item \emph{Gap}: $(1830-1723{,}3)/1830=5{,}83\,\%$.
  \item Con $M=10^{4}$: costes efectivos $12{,}04$, $16{,}015$ y $14{,}025$~€/t; $z_{\mathrm{LP}}=80\cdot 12{,}04+20\cdot 14{,}025=1243{,}7$~€ y el gap sube al $32{,}0\,\%$: la relajación casi no paga los costes fijos ($y_1=0{,}008$).
\end{enumerate}

\noindent\textbf{(e)}
\begin{enumerate}[label=\arabic*.]
  \item «$y_2=1\Rightarrow x_1\le 60$»: $x_1\le 60+M(1-y_2)$. Como $x_1\le 80$ siempre, la menor constante válida es $M=80-60=20$: $x_1\le 80-20\,y_2$.
  \item Solo cambian los patrones con $y_2=1$ y $x_1>60$. $\{1,2\}$: $x_1=60$, $x_2=40$: $550+720+640=1910$~€. $\{1,2,3\}$ no cambia ($x_1=55$).
  \item Nuevo óptimo: $\{2,3\}$ con $x=(0;\ 40;\ 60)$~t y $1880$~€.
\end{enumerate}
\resultado (b) $M=(80;\ 50;\ 60)$~t. (c) Hornos 1 y~2, $x=(80;\ 20;\ 0)$~t, $1830$~€. (d) $z_{\mathrm{LP}}=1723$~€ y gap del $5{,}83\,\%$ con las $M$ mínimas; $1244$~€ y $32{,}0\,\%$ con $M=10^{4}$: una $M$ grande no cambia el óptimo entero, pero destruye la cota y agranda el árbol. (e) $x_1\le 80-20y_2$; óptimo $\{2,3\}$, $x=(0;\ 40;\ 60)$~t, $1880$~€.
\comprobacion la cota es coherente: $1723{,}3\le 1830$ (proposición~2.10). En el nuevo óptimo se cumplen $L_2=20\le 40\le 50$ y $L_3=25\le 60\le 60$, y la igualdad $40+60=100$. Un PLEM con HiGHS da los mismos valores.

[Criterios: (a) 1~p por la función objetivo, 1~p por el enlace y 1~p por el lote mínimo y la igualdad. (b) 1~p por los valores y 1~p por la justificación. (c) 0{,}5~p por los patrones infactibles y 2{,}5~p por los cuatro patrones factibles y la conclusión. (d) 1~p por cada relajación con su gap. (e) 1~p por la restricción con $M=20$ y 1~p por el nuevo óptimo.]

Repasar: M24, §2.1, §2.3 y §2.4.
\end{solucion}

% ------------------------------------------------------------- B.2
\begin{solucion}{eje:bo}
\datos núcleo exponencial cuadrático con $\sigma_f=\ell=1$; $X=(0,2)$, $y=(0{,}5;\ -0{,}8)$; $\fmejor=-0{,}8$.
\planteamiento posterior del proceso gaussiano (teorema~3.5 de M25), mejora esperada en forma cerrada (teorema~4.4), probabilidad de mejora (definición~4.2) y cota inferior de confianza (definición~4.7).
\pasos

\noindent\textbf{(a)}
\begin{enumerate}[label=\arabic*.]
  \item $\rho=k(0,2)=e^{-2}=0{,}1353$; $K=\begin{psmallmatrix}1&\rho\\\rho&1\end{psmallmatrix}$, $\det K=1-\rho^{2}=0{,}9817$ y $K^{-1}=\frac{1}{0{,}9817}\begin{psmallmatrix}1&-0{,}1353\\-0{,}1353&1\end{psmallmatrix}=\begin{psmallmatrix}1{,}0187&-0{,}1379\\-0{,}1379&1{,}0187\end{psmallmatrix}$.
  \item $\alpha=K^{-1}y=\frac{1}{0{,}9817}(0{,}5+0{,}1353\cdot 0{,}8;\ -0{,}8-0{,}1353\cdot 0{,}5)=(0{,}6196;\ -0{,}8839)$.
  \item En $x=1$: $k_\ast=(e^{-1/2},e^{-1/2})=(0{,}6065;\ 0{,}6065)$. $\mu(1)=0{,}6065\,(0{,}6196-0{,}8839)=-0{,}1603$.
  \item Como $k_\ast\propto(1,1)$ y $K^{-1}(1,1)\T=(1,1)\T/(1+\rho)$: $\sigma^{2}(1)=1-2\cdot 0{,}6065^{2}/1{,}1353=1-0{,}6481=0{,}3519$; $\sigma(1)=0{,}5932$.
\end{enumerate}

\noindent\textbf{(b)}
\begin{enumerate}[label=\arabic*.]
  \item $k_\ast(3)=(e^{-9/2},e^{-1/2})=(0{,}01111;\ 0{,}6065)$.
  \item $\mu(3)=0{,}01111\cdot 0{,}6196+0{,}6065\cdot(-0{,}8839)=0{,}00688-0{,}5361=-0{,}5292$.
  \item $k_\ast\T K^{-1}k_\ast=(0{,}01111^{2}-2\cdot 0{,}1353\cdot 0{,}01111\cdot 0{,}6065+0{,}6065^{2})/0{,}9817=0{,}3661/0{,}9817=0{,}3730$; $\sigma^{2}(3)=0{,}6270$, $\sigma(3)=0{,}7918$.
\end{enumerate}

\noindent\textbf{(c)}
\begin{enumerate}[label=\arabic*.]
  \item $x=1$: $z=(-0{,}8+0{,}1603)/0{,}5932=-1{,}078$; $\Phi(z)=0{,}1404$, $\varphi(z)=0{,}2231$ (interpolando en el \cref{tab:normal}). $\PI=0{,}140$; $\EI=-0{,}6397\cdot 0{,}1404+0{,}5932\cdot 0{,}2231=-0{,}0898+0{,}1324=0{,}0425$.
  \item $x=3$: $z=(-0{,}8+0{,}5292)/0{,}7918=-0{,}3420$; $\Phi(z)=0{,}3662$, $\varphi(z)=0{,}3763$. $\PI=0{,}366$; $\EI=-0{,}2708\cdot 0{,}3662+0{,}7918\cdot 0{,}3763=-0{,}0992+0{,}2980=0{,}199$.
  \item $\LCB(1)=-0{,}1603-2\cdot 0{,}5932=-1{,}347$; $\LCB(3)=-0{,}5292-2\cdot 0{,}7918=-2{,}113$.
  \item Los tres criterios eligen $x=3$. No hace falta calcular para saberlo: $x=3$ tiene a la vez menor media y mayor desviación, y la EI decrece con $\mu$ y crece con $\sigma$ (proposición~4.5 de M25). El punto 3 combina explotación (la tendencia apunta a valores más bajos más allá del mejor ensayo) y exploración (está lejos de los datos); $x=1$, entre los dos ensayos, está bien acotado y es peor que $\fmejor$ con probabilidad alta.
\end{enumerate}

\noindent\textbf{(d)}
\begin{enumerate}[label=\arabic*.]
  \item $K_y=K+0{,}1\,I$, $\det K_y=1{,}1^{2}-0{,}1353^{2}=1{,}1917$; $k_\ast(2)=(0{,}1353;\ 1)$.
  \item $k_\ast\T K_y^{-1}k_\ast=(1{,}1\cdot 0{,}01832-2\cdot 0{,}1353\cdot 0{,}1353+1{,}1)/1{,}1917=1{,}0835/1{,}1917=0{,}9092$; $\sigma^{2}_{f}(2)=0{,}0908$.
  \item Varianza de una nueva observación: $0{,}0908+0{,}1=0{,}191$.
  \item En la EI se usa la latente, $0{,}0908$: la mejora se refiere al valor verdadero $f(x)$, no a una medida ruidosa. Con la predictiva, la incertidumbre de un punto ensayado nunca baja de $\sigma_\varepsilon$ y el algoritmo repite ensayos (observación~4.8 de M25).
\end{enumerate}

\noindent\textbf{(e)}
\begin{enumerate}[label=\arabic*.]
  \item Con $\ell=0{,}3$: $k(0,1)=k(2,1)=e^{-0{,}5/0{,}09}=e^{-5{,}556}=0{,}00387$ y $k(0,2)=e^{-22{,}2}\approx 0$, de modo que $K\approx I$ y $\alpha\approx y$.
  \item $\mu(1)\approx 0{,}00387\,(0{,}5-0{,}8)=-0{,}0012\approx 0$; $\sigma^{2}(1)\approx 1-2\cdot 0{,}00387^{2}\approx 1$.
  \item $z\approx-0{,}8$: $\EI(1)\approx-0{,}8\cdot 0{,}2119+0{,}2897=0{,}120$. En $x=3$ sale prácticamente lo mismo ($0{,}121$).
  \item Con una longitud de escala corta los datos no informan entre ensayos: toda la zona no ensayada vuelve a la previa y la EI es casi uniforme lejos de los datos, de modo que el algoritmo explora casi al azar. Por eso $\ell$ se estima por verosimilitud marginal (M25, §3.4) y se vigila que no colapse.
\end{enumerate}
\resultado (a) $\alpha=(0{,}620;\ -0{,}884)$, $\mu(1)=-0{,}160$, $\sigma^{2}(1)=0{,}352$. (b) $\mu(3)=-0{,}529$, $\sigma^{2}(3)=0{,}627$. (c) $\PI=(0{,}140;\ 0{,}366)$, $\EI=(0{,}0425;\ 0{,}199)$, $\LCB=(-1{,}35;\ -2{,}11)$; los tres criterios eligen $x=3$. (d) $0{,}0908$ (latente) y $0{,}191$ (observación); se usa la latente. (e) $\mu(1)\approx 0$, $\sigma^{2}(1)\approx 1$, $\EI(1)\approx 0{,}120$.
\comprobacion la varianza nunca supera la previa ($0{,}352<1$, $0{,}627<1$). Ambos candidatos distan una unidad de su ensayo más próximo, pero $x=1$ tiene \emph{dos} ensayos a esa distancia y $x=3$ solo uno, de ahí que su varianza sea menor. Una búsqueda en rejilla de la EI con $\ell=1$ da el máximo en $x\approx 2{,}95$ con $\EI=0{,}199$.

[Criterios: (a) 1~p por $K^{-1}$ y $\alpha$, 1~p por $\mu(1)$ y 1~p por $\sigma^{2}(1)$. (b) 1~p por cada cifra. (c) 1~p por las EI, 1~p por PI y LCB y 1~p por la explicación. (d) 1~p por los cálculos y 1~p por la elección razonada. (e) 1~p por las cifras y 1~p por la interpretación.]

Repasar: M25, §3.3, §3.4, §4.2, §4.3, §4.4 y §4.5.
\end{solucion}

% ------------------------------------------------------------- B.3
\begin{solucion}{eje:mpc}
\datos $a=0{,}9$, $b=0{,}5$; $Q=1$, $R=0{,}5$, $P_f=1$, $N=2$, $x_0=2$; en~(c), $q=2$, $r=1$, $\gamma=0{,}95$.
\planteamiento (a) recursión de Riccati del teorema~3.2 de M26; (b) forma condensada y QP de la proposición~3.6, resuelto con las condiciones de Karush--Kuhn--Tucker como en el ejemplo~3.7; (c) ecuación de Bellman (M26, §3.1 y §4.1) con una función de valor afín ensayada.
\pasos

\noindent\textbf{(a)}
\begin{enumerate}[label=\arabic*.]
  \item $P_2=1$. $K_1=bP_2a/(R+b^{2}P_2)=0{,}45/0{,}75=0{,}6$; $P_1=Q+a^{2}P_2-aP_2bK_1=1+0{,}81-0{,}27=1{,}54$.
  \item $K_0=bP_1a/(R+b^{2}P_1)=0{,}693/0{,}885=0{,}7831$; $P_0=1+0{,}81\cdot 1{,}54-0{,}9\cdot 1{,}54\cdot 0{,}5\cdot 0{,}7831=1+1{,}2474-0{,}5426=1{,}7047$.
  \item $u_0=-0{,}7831\cdot 2=-1{,}566$; $x_1=1{,}8-0{,}7831=1{,}017$; $u_1=-0{,}6\cdot 1{,}017=-0{,}6102$; $x_2=0{,}9153-0{,}3051=0{,}6102$.
  \item Coste: $4+0{,}5\cdot 2{,}453+1{,}034+0{,}5\cdot 0{,}3723+0{,}3723=6{,}819=P_0x_0^{2}=1{,}7047\cdot 4$.
\end{enumerate}

\noindent\textbf{(b)}
\begin{enumerate}[label=\arabic*.]
  \item $\Phi=(a,a^{2})\T=(0{,}9;\ 0{,}81)\T$, $\Gamma=\begin{psmallmatrix}b&0\\ab&b\end{psmallmatrix}=\begin{psmallmatrix}0{,}5&0\\0{,}45&0{,}5\end{psmallmatrix}$, $\bar Q=\operatorname{diag}(1,1)$, $\bar R=\operatorname{diag}(0{,}5;\ 0{,}5)$.
  \item $H=\Gamma\T\bar Q\Gamma+\bar R=\begin{psmallmatrix}0{,}9525&0{,}225\\0{,}225&0{,}75\end{psmallmatrix}$; $F=\Gamma\T\bar Q\Phi=(0{,}8145;\ 0{,}405)$; constante $x_0^{2}(1+0{,}81+0{,}6561)=9{,}864$.
  \item El óptimo libre tiene $u_0=-1{,}566<-1$: viola la cota. Se ensaya $u_0=-1$ activo y se minimiza en $u_1$: $\partial J/\partial u_1=2(0{,}225u_0+0{,}75u_1)+2\cdot 2\cdot 0{,}405=0$ da $u_1=(0{,}225-0{,}81)/0{,}75=-0{,}78$, dentro de la cota.
  \item $J=7{,}103$ (frente a $6{,}819$ sin restricción). En ese punto $\partial J/\partial u_0=2(0{,}9525\cdot(-1)+0{,}225\cdot(-0{,}78))+2\cdot 2\cdot 0{,}8145=1{,}002>0$: subir $u_0$ empeora y bajarlo no está permitido; el multiplicador de $u_0\ge-1$ es $1{,}002$.
  \item Se aplica solo $u_0=-1$~kW; en el instante siguiente se mide $x_1$ y se vuelve a resolver (horizonte deslizante, M26, §3.2).
\end{enumerate}

\noindent\textbf{(c)}
\begin{enumerate}[label=\arabic*.]
  \item Bellman: $V(z)=\min_u\bigl\{qz+ru^{2}+\gamma V(az+bu)\bigr\}$.
  \item Con $V(z)=pz+s$: $pz+s=qz+\gamma paz+\gamma s+\min_u\{ru^{2}+\gamma pbu\}$.
  \item El mínimo en $u$ no depende de $z$: $u^{\ast}=-\gamma pb/(2r)$, con valor $-(\gamma pb)^{2}/(4r)$.
  \item Identificando coeficientes: $p=q+\gamma ap$, luego $p=q/(1-\gamma a)=2/0{,}145=13{,}79$; y $s(1-\gamma)=-(\gamma pb)^{2}/(4r)$.
  \item $\gamma pb=0{,}95\cdot 13{,}79\cdot 0{,}5=6{,}552$; $u^{\ast}=-3{,}276$~kW; $s=-6{,}552^{2}/(4\cdot 0{,}05)=-214{,}6$.
\end{enumerate}

\noindent\textbf{(d)}
\begin{enumerate}[label=\arabic*.]
  \item Si $V(z)=pz+s$, el término en $z$ de la ecuación de Bellman no depende de $u$, de modo que la minimización en $u\in[-2,2]$ sigue sin depender de $z$ y la identificación $p=q+\gamma ap$ no cambia: la forma afín se conserva con la misma $p$.
  \item La función $ru^{2}+\gamma pbu$ es convexa con mínimo libre en $-3{,}276$, fuera del intervalo: $u^{\ast}=-2$~kW (saturada).
  \item $s(1-\gamma)=r\cdot 4+6{,}552\cdot(-2)=-9{,}103$, $s=-182{,}1$.
\end{enumerate}
\resultado (a) $K_1=0{,}600$, $P_1=1{,}54$, $K_0=0{,}783$, $P_0=1{,}70$; $u=(-1{,}57;\ -0{,}610)$~kW y $J=6{,}82=P_0x_0^{2}$. (b) $u=(-1;\ -0{,}780)$~kW, $J=7{,}10$, multiplicador $1{,}00$; se aplica $u_0=-1$~kW. (c) $V(z)=13{,}8\,z-215$ y $u^{\ast}=-3{,}28$~kW, constante. (d) $u^{\ast}=-2$~kW y $V(z)=13{,}8\,z-182$.
\comprobacion en~(c), simulando la política constante desde $z=1$ durante 3000 pasos se obtiene $J=-200{,}83=13{,}79-214{,}63$, y perturbar $u$ en $\pm0{,}1$ empeora el coste en $0{,}2$ en ambos sentidos. La condición $\gamma a=0{,}855<1$ garantiza que la serie converge. La acción no depende del estado porque el coste marginal de tener una unidad más de $x$ es el mismo en cualquier estado: es lo que dice la pendiente constante $p$.

[Criterios: (a) 1~p por $K_1,P_1$, 1~p por $K_0,P_0$ y 1~p por la secuencia y la comprobación. (b) 1~p por $H$ y $F$, 1~p por la solución con la cota activa y 1~p por el multiplicador y el principio de horizonte deslizante. (c) 1~p por la ecuación de Bellman, 1~p por $u^{\ast}$, 1~p por $p$ y 1~p por $s$. (d) 1~p por el argumento y 1~p por $u^{\ast}$ y $s$.]

Repasar: M26, §3.1, §3.2, §3.3 y §4.1.
\end{solucion}

% ------------------------------------------------------------- B.4
\begin{solucion}{eje:simp}
\datos $k_0=1000$~N/mm, $p=3$, $x=(1;\ 0{,}8;\ 0{,}5)$, $F=(0,100,100)$~N.
\planteamiento la barra en serie es isostática; el gradiente de la flexibilidad se obtiene con el corolario~2.5 de M27 y el de $J=e_2\T U$ con el teorema~2.4 (método adjunto); la actualización es la de la definición~2.7, y con $p=1$ el problema es convexo (teorema~2.3), de modo que las condiciones de Karush--Kuhn--Tucker son suficientes (M24, §4.2).
\pasos

\noindent\textbf{(a)}
\begin{enumerate}[label=\arabic*.]
  \item $k=(1000;\ 512;\ 125)$~N/mm. $K=\begin{psmallmatrix}k_1+k_2&-k_2&0\\-k_2&k_2+k_3&-k_3\\0&-k_3&k_3\end{psmallmatrix}=\begin{psmallmatrix}1512&-512&0\\-512&637&-125\\0&-125&125\end{psmallmatrix}$.
  \item Esfuerzos axiles (equilibrio de cada tramo): $N=(200;\ 200;\ 100)$~N. Alargamientos $\delta_e=N_e/k_e=(0{,}2;\ 0{,}3906;\ 0{,}8)$~mm.
  \item $U=(0{,}2;\ 0{,}5906;\ 1{,}3906)$~mm; $c=100\cdot 0{,}5906+100\cdot 1{,}3906=198{,}1$~N·mm.
\end{enumerate}

\noindent\textbf{(b)}
\begin{enumerate}[label=\arabic*.]
  \item $g=F\T U$ no depende de $x$ y $(\partial g/\partial U)\T=F$; el adjunto resuelve $K\lambda=F$, luego $\lambda=U$.
  \item $\partial c/\partial x_e=-U\T(\partial K/\partial x_e)U=-p\,x_e^{p-1}k_0\,\delta_e^{2}$: $-3\cdot 1\cdot 1000\cdot 0{,}04=-120$; $-3\cdot 0{,}64\cdot 1000\cdot 0{,}1526=-293{,}0$; $-3\cdot 0{,}25\cdot 1000\cdot 0{,}64=-480$~N·mm.
  \item Comprobación directa: $c=\sum_eN_e^{2}/(k_0x_e^{3})$, con $N$ fijo; $\partial c/\partial x_3=-3N_3^{2}/(k_0x_3^{4})=-3\cdot 10^{4}/(1000\cdot 0{,}0625)=-480$.
\end{enumerate}

\noindent\textbf{(c)}
\begin{enumerate}[label=\arabic*.]
  \item Adjunto: $K\lambda=e_2$, es decir, el desplazamiento producido por una carga unidad en el nodo~2: $\lambda=(1/k_1;\ 1/k_1+1/k_2;\ 1/k_1+1/k_2)=(0{,}001;\ 0{,}002953;\ 0{,}002953)$~mm/N, con alargamientos $\delta^{\lambda}=(0{,}001;\ 0{,}001953;\ 0)$.
  \item $\partial J/\partial x_e=-\lambda\T(\partial K/\partial x_e)U=-p\,x_e^{p-1}k_0\,\delta^{\lambda}_e\,\delta_e$: $-3\cdot 1000\cdot 0{,}001\cdot 0{,}2=-0{,}6$; $-3\cdot 0{,}64\cdot 1000\cdot 0{,}001953\cdot 0{,}3906=-1{,}465$; $0$~mm.
  \item El elemento~3 está aguas abajo del nodo~2: su rigidez no cambia $u_2$, porque la barra es isostática. Comprobación: $u_2=200/k_1+200/k_2$ no depende de $x_3$ y $\partial u_2/\partial x_1=-200\cdot 3000/10^{6}=-0{,}6$.
\end{enumerate}

\noindent\textbf{(d)}
\begin{enumerate}[label=\arabic*.]
  \item $B_e=(-\partial c/\partial x_e)/\Lambda$ con $\partial V/\partial x_e=1$. Con $\Lambda=300$: $B=(0{,}4;\ 0{,}9766;\ 1{,}6)$, $\sqrt B=(0{,}6325;\ 0{,}9882;\ 1{,}2649)$.
  \item $x\sqrt B=(0{,}6325;\ 0{,}7906;\ 0{,}6325)$. Los límites de movimiento son $[0{,}8;\,1]$, $[0{,}6;\,1]$ y $[0{,}3;\,0{,}7]$, de modo que $x^{\text{nuevo}}=(0{,}8;\ 0{,}7906;\ 0{,}6325)$, con suma $2{,}223<2{,}3$.
  \item Falta material: $\Lambda$ (el precio del material) es demasiado alto y debe bajar (proposición~2.8: el volumen decrece con $\Lambda$).
  \item Con $\Lambda=270$: $\sqrt{120/270}=0{,}6667$ (se recorta a $0{,}8$), $\sqrt{292{,}97/270}=1{,}0417$ y $\sqrt{480/270}=1{,}3333$; $x^{\text{nuevo}}=(0{,}8;\ 0{,}8333;\ 0{,}6667)$, suma $2{,}3$.
\end{enumerate}

\noindent\textbf{(e)}
\begin{enumerate}[label=\arabic*.]
  \item Con $p=1$: $c(x)=\sum_eN_e^{2}/(k_0x_e)$, convexa en $x>0$. Lagrangiana $c+\Lambda(\sum_ex_e-2{,}3)$; estacionariedad: $N_e^{2}/(k_0x_e^{2})=\Lambda$, luego $x_e\propto N_e$.
  \item $x_e=2{,}3\,N_e/\sum N=2{,}3\cdot(200,200,100)/500=(0{,}92;\ 0{,}92;\ 0{,}46)$, todas $\le 1$; $\Lambda=(500/2{,}3)^{2}/1000=47{,}26$~N·mm por unidad de volumen.
  \item $c^{\ast}=\frac{1}{1000}\bigl(\frac{4\cdot 10^{4}}{0{,}92}\cdot 2+\frac{10^{4}}{0{,}46}\bigr)=108{,}7$~N·mm, frente a $c(x_0)=\frac{1}{1000}\bigl(\frac{4\cdot 10^{4}}{1}+\frac{4\cdot 10^{4}}{0{,}8}+\frac{10^{4}}{0{,}5}\bigr)=110{,}0$~N·mm.
\end{enumerate}
\resultado (a) Desplazamientos $U=(0{,}200;\ 0{,}591;\ 1{,}39)$~mm y flexibilidad $c=198$~N·mm. (b) Gradiente $\nabla c=(-120;\ -293;\ -480)$~N·mm. (c) Gradiente $\nabla J=(-0{,}600;\ -1{,}46;\ 0)$~mm. (d) Con $\Lambda=300$ resulta $x=(0{,}800;\ 0{,}791;\ 0{,}633)$ y un volumen de $2{,}22$, de modo que hay que bajar $\Lambda$; con $\Lambda=270$ resulta $x=(0{,}800;\ 0{,}833;\ 0{,}667)$ y un volumen de $2{,}30$. (e) $x^{\ast}=(0{,}920;\ 0{,}920;\ 0{,}460)$ y $c^{\ast}=109$~N·mm, un $1{,}2\,\%$ menos que el diseño inicial con $p=1$.
\comprobacion las diferencias finitas con paso $10^{-6}$ reproducen los tres gradientes de~(b) y de~(c) con cinco cifras. En~(e), en el óptimo $N_e^{2}/(k_0x_e^{2})$ vale $47{,}26$ en los tres elementos, como exige la estacionariedad, y el signo negativo de todas las sensibilidades de~(b) confirma que añadir material nunca aumenta la flexibilidad.

[Criterios: (a) 1~p por $K$ y 1~p por $U$ y $c$. (b) 1~p por el argumento autoadjunto, 1~p por las tres derivadas y 1~p por la comprobación. (c) 1~p por el problema adjunto, 1~p por las derivadas y 1~p por la interpretación. (d) 1~p por $x^{\text{nuevo}}$ y el sentido de $\Lambda$ y 1~p por la verificación con $\Lambda=270$. (e) 1~p por las condiciones KKT y 1~p por el óptimo y la comparación.]

Repasar: M27, §2.1, §2.2, §2.3 y §2.4; M24, §4.2.
\end{solucion}

% ------------------------------------------------------------- B.5
\begin{solucion}{eje:valor}
\datos $n=16$, $\bar d=1200$~€/semana, $s_d=1600$~€/semana; $\sigma_A=\sigma_B=3000$~€, $r=0{,}86$; percentiles \emph{bootstrap} 560 y 1890~€/semana; $I=60\,000$~€, $O=12\,000$~€/año, $N=5$, $\rho=0{,}08$.
\planteamiento intervalo $t$ pareado y ventaja del emparejamiento (proposición~5.3 de M28), intervalos \emph{bootstrap} (definición~5.4), VAN con anualidad (proposición~5.2) y tamaño de un piloto (proposición~5.5).
\pasos

\noindent\textbf{(a)}
\begin{enumerate}[label=\arabic*.]
  \item Error típico $s_d/\sqrt n=1600/4=400$; semiamplitud $2{,}131\cdot 400=852{,}4$.
  \item Intervalo: $[347{,}6;\ 2052{,}4]$~€/semana, es decir, $[18\,075;\ 106\,725]$~€/año.
  \item Hipótesis: diferencias independientes entre semanas (o autocorrelación despreciable, que debe comprobarse) y aproximadamente normales, o $n$ suficiente; el gemelo debe estar validado, porque el intervalo mide la variabilidad entre semanas y no el error del modelo.
\end{enumerate}

\noindent\textbf{(b)}
\begin{enumerate}[label=\arabic*.]
  \item $\sigma_D=\sqrt{3000^{2}+3000^{2}-2\cdot 0{,}86\cdot 3000^{2}}=\sqrt{2{,}52\cdot 10^{6}}=1587$~€, coherente con $s_d$.
  \item Sin emparejar: $\sqrt{2\cdot 3000^{2}}=4243$~€.
  \item Cociente de varianzas: $1/(1-r)=7{,}14$: harían falta unas siete veces más semanas.
\end{enumerate}

\noindent\textbf{(c)}
\begin{enumerate}[label=\arabic*.]
  \item Percentil: $[560;\ 1890]$~€/semana.
  \item Básico: $[2\bar d-1890;\ 2\bar d-560]=[510;\ 1840]$~€/semana.
  \item La distribución de las diferencias es asimétrica: la mayoría de semanas el ahorro es moderado y, de vez en cuando, el optimizador evita una avería cara. El \emph{bootstrap} recoge esa asimetría sin suponer normalidad; ambos métodos suponen semanas independientes.
\end{enumerate}

\noindent\textbf{(d)}
\begin{enumerate}[label=\arabic*.]
  \item $a_{5;\,0{,}08}=(1-1{,}08^{-5})/0{,}08=3{,}9927$.
  \item Con la media: $b=52\cdot 1200=62\,400$~€/año; $\VAN=(62\,400-12\,000)\cdot 3{,}9927-60\,000=141\,233$~€.
  \item Con el extremo inferior: $b=18\,075$; $\VAN=(18\,075-12\,000)\cdot 3{,}9927-60\,000=-35\,744$~€.
  \item Umbral: $\VAN=0\iff b=O+I/a=12\,000+15\,027=27\,027$~€/año, es decir, $519{,}8$~€/semana.
  \item Como $347{,}6<519{,}8$, el intervalo del 95\,\% incluye mejoras que no recuperan la inversión: no puede afirmarse con esa confianza que el proyecto sea rentable, aunque el VAN esperado sea claramente positivo. (El contraste unilateral de $H_0\colon\delta\le 519{,}8$ da $t=(1200-519{,}8)/400=1{,}70<1{,}753$.)
\end{enumerate}

\noindent\textbf{(e)}
\begin{enumerate}[label=\arabic*.]
  \item $z_{0{,}975}+z_{0{,}80}=1{,}960+0{,}8416=2{,}802$.
  \item $n\approx(2{,}802\cdot 1600/519{,}8)^{2}=(8{,}624)^{2}=74{,}4$: unas 75 semanas.
\end{enumerate}
\resultado (a) $[348;\ 2050]$~€/semana, es decir, $[18\,100;\ 107\,000]$~€/año. (b) $\sigma_D=1590$~€ frente a $4240$~€; unas $7{,}1$ veces más semanas. (c) Percentil $[560;\ 1890]$ y básico $[510;\ 1840]$~€/semana. (d) $\VAN=141\,000$~€ con la media y $-35\,700$~€ con el extremo inferior; la rentabilidad nula exige $520$~€/semana; no puede afirmarse con un 95\,\% de confianza. (e) Unas 75 semanas.
\comprobacion la recuperación con la mejora media es $60\,000/50\,400=1{,}19$ años y la tasa interna de rentabilidad, del $79\,\%$; con $\delta=1200$ bastarían $n\approx 14$ semanas, lo que explica que 16 semanas detecten la mejora pero no basten para demostrar la rentabilidad. Las conclusiones serían las de M28, §5.3: completar con más líneas, un periodo en sombra o indicadores con menos ruido.

[Criterios: (a) 1~p por la semiamplitud, 1~p por el intervalo en las dos escalas y 1~p por las hipótesis. (b) 1~p por las dos desviaciones y 1~p por el factor. (c) 1~p por los dos intervalos y 1~p por la explicación. (d) 1~p por los dos VAN, 1~p por el umbral y 1~p por la conclusión. (e) 2~p, o 1~p si se usa $\delta=1200$ sin justificarlo.]

Repasar: M28, §5.1, §5.2 y §5.3.
\end{solucion}
